\section{Implementação do sistema}
\label{sec:implementacao-real}
Esta versão do protótipo foi implementada em Java e inspecionada em 29 de setembro de 2026. O projeto é uma implementação paralela do mesmo recorte experimental, não uma comparação entre linguagens. O código-fonte e sua configuração estão no repositório \url{https://github.com/ilann47/TCC-2026}, no diretório \texttt{implementacao-java}. As verificações funcionais descritas nesta seção não constituem a coleta quantitativa definitiva prevista no Capítulo~\ref{cap:metodo}.

\subsection{Tecnologias e módulos}
O projeto utiliza Java com bytecode alvo 21, Spring Boot 3.5.16, Maven Wrapper, driver PostgreSQL, Spring JDBC, Flyway, cliente Apache Kafka 3.9.1 e springdoc-openapi 2.9.1 para a documentação das APIs. A implantação descrita em \texttt{docker-compose.yml} utiliza PostgreSQL 16 e Kafka 3.9.1. O Maven compila cinco módulos: \texttt{audit-core}, \texttt{audit-storage}, \texttt{sync-api}, \texttt{async-api} e \texttt{consumer}. Os dois primeiros são bibliotecas compartilhadas; os três últimos são processos executáveis independentes.

\begin{table}[htbp]
\centering\small
\caption{Módulos Java e responsabilidades efetivas}\label{tab:modulos-implementados}
\begin{tabular}{p{0.24\textwidth}p{0.65\textwidth}}
\hline
Módulo & Responsabilidade\\ \hline
\texttt{audit-core} & Contrato \texttt{AuditEvent}, serialização determinística, SHA-256, publicação Kafka e marcos de observabilidade.\\
\texttt{audit-storage} & Migração Flyway V1, transação JDBC, consulta, comparação de hash e persistência idempotente.\\
\texttt{sync-api} & \texttt{POST /audit}, \texttt{GET /audit/\{event\_id\}}, \texttt{/live} e \texttt{/health}; grava antes da resposta 201.\\
\texttt{async-api} & \texttt{POST /audit}, \texttt{/live} e \texttt{/health}; publica no Kafka e responde 202 depois do ACK. Não acessa o banco diretamente.\\
\texttt{consumer} & Consome o tópico, persiste, aplica retry, publica na DLQ quando necessário e confirma o offset depois da resolução.\\ \hline
\end{tabular}
\fonte{Elaborado pelo autor a partir do código Java (2026).}
\end{table}

O Compose inicia os três processos executáveis, PostgreSQL e Kafka. A API síncrona expõe a porta local 18002 e a assíncrona, 18003; PostgreSQL e Kafka expõem 15433 e 19093 somente em \textit{loopback}. O consumidor não expõe HTTP. Os serviços auxiliares \texttt{kafka-init} e \texttt{kafka-volume-init} preparam os tópicos e as permissões antes do uso. O ambiente possui um broker e uma partição por tópico; não representa alta disponibilidade. Os limites configurados são 1 CPU e 512 MiB para cada processo Java, 1 CPU e 512 MiB para PostgreSQL, e 1 CPU e 1 GiB para Kafka. Esses valores delimitam a configuração de desenvolvimento, sem constituir calibração científica.

A API síncrona hospeda uma única página Swagger UI na porta local 18002. Seu seletor carrega separadamente os contratos OpenAPI das variantes síncrona e assíncrona; a segunda API expõe o documento JSON, mas não outra interface Swagger. A configuração de CORS da API assíncrona permite à origem local da página consultar o contrato e executar requisições de teste. Essa ferramenta auxilia a inspeção manual das respostas 201 e 202, sem constituir frontend próprio, gerador de carga ou instrumento de medição do experimento.

\subsection{Comportamento implementado}
O contrato \texttt{AuditEvent} é compartilhado pelas duas APIs. Campos de texto obrigatórios respeitam os comprimentos definidos no DTO; \texttt{event\_id}, \texttt{occurred\_at} e \texttt{payload} recebem valores padrão apenas quando omitidos no HTTP. O consumidor exige UUID e instante explícitos na mensagem recebida para não regenerar identidade em uma reentrega. A serialização canônica e o SHA-256 são comuns às variantes.

A API síncrona chama \texttt{AuditRepository.persist}, que executa \texttt{INSERT ON CONFLICT DO NOTHING} em transação \texttt{READ COMMITTED}. Uma segunda submissão com UUID e hash iguais preserva uma linha e retorna \texttt{duplicate}; hash diferente provoca conflito 409 e mantém o conteúdo original. O retorno 201 ocorre após o término da transação. \texttt{persisted\_at} é o horário atribuído durante a inserção, enquanto o marco \texttt{commit\_completed} é emitido depois do retorno da transação.

A API assíncrona usa \texttt{KafkaPublisher.publish}, que aguarda a conclusão de \texttt{send(...).get(...)} dentro de um prazo configurado. \texttt{acks=all} e a idempotência do produtor estão habilitados. O retorno 202 confirma a publicação no broker configurado, mas não confirma a gravação no PostgreSQL. Uma falha ou timeout de publicação retorna 503; em alguns casos a aceitação permanece incerta, exigindo conciliação pelo UUID antes de classificar o evento como perdido.

O consumidor usa commit manual de offsets. Após persistência ou reconhecimento idempotente, confirma a posição de origem. Mensagem inválida, conflito ou esgotamento das tentativas de banco seguem para \texttt{audit-events-dlq}; o offset de origem só é confirmado depois do ACK da DLQ. Se a publicação na DLQ ou o commit do offset falhar, a sessão é retomada sem confirmação segura da posição. A DLQ pode conter cópias repetidas após falha entre seu ACK e o commit de origem. Não existe reprocessador automático da DLQ: um evento aceito com 202 durante indisponibilidade prolongada do banco pode terminar em quarentena e não será persistido automaticamente quando o banco voltar.

\subsection{Verificação funcional e limites da evidência}
\label{sec:validacao-funcional}
O Maven Wrapper executa build e testes unitários dos módulos Java. O fluxo de integração contínua publicado no repositório executou \texttt{mvnw verify}, iniciou PostgreSQL e Kafka em contêineres, verificou o caminho normal de registro síncrono e assíncrono e conferiu o seletor da documentação e a política CORS necessária ao teste manual. Um \textit{smoke test} de C4 interrompeu o PostgreSQL e verificou a resposta 503 da API síncrona, a aceitação 202 da API assíncrona enquanto Kafka permanecia disponível e a retomada da saúde do banco. O registro de execução do fluxo está disponível em \url{https://github.com/ilann47/TCC-2026/actions/runs/36647710056}.

Essas verificações demonstram apenas os comportamentos observados no ambiente de CI. O \textit{smoke test} C4 não prova que cada evento aceito durante a falha seja persistido depois da recuperação; tampouco mede latência, throughput ou tempo de drenagem de maneira comparativa. Em outubro de 2026, foram acrescentados scripts k6, um executor isolado, um analisador independente Java e um comparador de exportações do modelo comum. Esses instrumentos possuem testes próprios e preservam evidências locais. Sua validação funcional e o piloto são discutidos no Capítulo~\ref{cap:resultados}, sem substituírem a calibração, o congelamento dos parâmetros e as repetições da campanha definitiva.

\subsection{Rastreabilidade da entrega}
\label{sec:evidencias}
As fontes de verificação desta versão são o código Java versionado, a migração SQL \texttt{V1\_\_create\_audit\_records.sql}, os testes dos módulos, o arquivo Compose e o fluxo de integração contínua. Os diagramas desta entrega foram refeitos para corresponder a essas fontes; modelos de interface e tabelas adicionais continuam identificados como propostas. O repositório preserva o histórico de commits e a especificação técnica versionável em \texttt{implementacao-java/docs/especificacao}. Resultados experimentais serão apresentados somente após protocolo congelado, execução das repetições e conciliação dos dados brutos.
